%%%%%%%%%%%%%%%%%%%%%%%%%%%%%%%%%%%%%%%%%%%%%%%%%%%%%%%%%%%%%%%%%%%%%%%%%%%%%%%%
%%%%%%%%%%%%%%%%%%%%%%%% ------------------------------ %%%%%%%%%%%%%%%%%%%%%%%%
%%%%%%%%%%%%%%%%%%%%%%%% | Formato escuela - INF PUCV | %%%%%%%%%%%%%%%%%%%%%%%%
%%%%%%%%%%%%%%%%%%%%%%%% | by Sebastian Garcia        | %%%%%%%%%%%%%%%%%%%%%%%%
%%%%%%%%%%%%%%%%%%%%%%%% | @sebaignacioo              | %%%%%%%%%%%%%%%%%%%%%%%%
%%%%%%%%%%%%%%%%%%%%%%%% |                            | %%%%%%%%%%%%%%%%%%%%%%%%
%%%%%%%%%%%%%%%%%%%%%%%% |            2023            | %%%%%%%%%%%%%%%%%%%%%%%%
%%%%%%%%%%%%%%%%%%%%%%%% ------------------------------ %%%%%%%%%%%%%%%%%%%%%%%%
%%%%%%%%%%%%%%%%%%%%%%%%%%%%%%%%%%%%%%%%%%%%%%%%%%%%%%%%%%%%%%%%%%%%%%%%%%%%%%%%
%
%  REORGANIZACIÓN (septiembre 2026): el contenido se separó en un
%  archivo .tex por sección, siguiendo el orden obligatorio del
%  documento "Indicaciones y fichas — Trabajo de Investigación 2026"
%  (FEP00.3.26) para el tema TI-05 (Serverless y computación en el
%  borde), empresa asignada Terabyte, cruzado además con las pautas
%  generales del curso (FEP00.1.26), que agregan "Marco teórico y
%  conceptual" y "Tendencias y evolución futura" y confirman el Anexo
%  de declaración de IA como obligatorio.
%
%  El preámbulo, la clase, los estilos y el formato de salida NO se
%  tocaron respecto de la plantilla original.
%
%  Desarrollo (bloque 5) y Entregables específicos (bloque 6) van cada
%  uno en un ÚNICO archivo .tex (con \subsection por subtema/
%  entregable), por preferencia del equipo — no se dividen en
%  subcarpetas ni archivos separados.
%
%  El contenido unificado original (previo a la primera reorganización)
%  queda respaldado, sin usarse, en _legacy/contenido.tex.bak.
%
%%%%%%%%%%%%%%%%%%%%%%%%%%%%%%%%%%%%%%%%%%%%%%%%%%%%%%%%%%%%%%%%%%%%%%%%%%%%%%%%

% Se define la clase y parametros principales
\documentclass[12pt, letterpaper, table, xcdraw]{inf-pucv}

% Usando paquete para importar archivos
\usepackage{import}
% Importando estiloss
\usepackage{styles}
\usepackage{subcaption}

% Modificar estilos de índice y listados
\renewcommand{\cftdotsep}{4}
\renewcommand{\cftfigpresnum}{\bfseries Figura }
\cftsetindents{figure}{0em}{5.5em}
\renewcommand{\cfttabpresnum}{\bfseries Tabla }
\cftsetindents{table}{0em}{5.5em}
% Se define la sangría de 1cm
\setlength{\parindent}{1cm}

\begin{document}

  %%%%%%%%%%%%%%%%%%%%%%%%%%%%%%%%%%%%%%%%%%%%%%%%%%%%%%%%%%%%%%%%%%%%%%%%%%%%%%
  % ÍNDICE DE BLOQUES — estructura obligatoria del informe TI-05, en
  % este orden. Cada bloque vive en su propio archivo dentro de
  % secciones/. Todas las secciones están vacías (solo encabezados y
  % comentarios guía), listas para redactar, EXCEPTO la Portada (bloque
  % 1), que ya tiene datos reales cargados.
  %
  %   Bloque  1  Portada y formalidades       secciones/00-portada.tex
  %   Bloque  2  Resumen ejecutivo            secciones/01-resumen-ejecutivo.tex
  %   Bloque  3  Introducción y contexto      secciones/02-introduccion.tex
  %   Bloque  4  Marco teórico y conceptual   secciones/03-marco-teorico.tex
  %   Bloque  5  Desarrollo y análisis        secciones/04-desarrollo.tex
  %              técnico (6 subtemas, 1 solo archivo)
  %   Bloque  6  Entregables específicos      secciones/05-entregables.tex
  %              (3 entregables, 1 solo archivo)
  %   Bloque  7  Vínculo técnico-económico    secciones/06-vinculo-tecnico-economico.tex
  %   Bloque  8  Tendencias y evolución       secciones/07-tendencias-futuras.tex
  %              futura
  %   Bloque  9  Conclusiones y               secciones/08-conclusiones.tex
  %              recomendaciones
  %   Bloque 10  Bibliografía                 secciones/09-bibliografia.tex
  %   Bloque 11  Anexo obligatorio:           secciones/10-anexo-declaracion-ia.tex
  %              declaración de uso de IA
  %   Bloque 12  Documento adjunto:           secciones/11-cuestionario.tex
  %              cuestionario de 30 preguntas
  %%%%%%%%%%%%%%%%%%%%%%%%%%%%%%%%%%%%%%%%%%%%%%%%%%%%%%%%%%%%%%%%%%%%%%%%%%%%%%

  %%%%%%%%%%%%%%%%%%%%%%%%%%%%%%%%%%%%%%%%%%%%%%%%%%%%%%%%%%%%%%%%%%%%%%%%%%%%%%%%
% BLOQUE 1 — PORTADA Y FORMALIDADES
% Estructura obligatoria del informe TI-05 (punto 1).
% Contiene: portada, indice general, lista de figuras y lista de tablas.
% Estos tres ultimos son formalidades de front-matter; por eso van sin
% numero de pagina visible (pagenumbering gobble). La numeracion arabiga
% del cuerpo del informe comienza en el Resumen ejecutivo (bloque 2).
%%%%%%%%%%%%%%%%%%%%%%%%%%%%%%%%%%%%%%%%%%%%%%%%%%%%%%%%%%%%%%%%%%%%%%%%%%%%%%%%

\thispagestyle{empty}
\pagenumbering{gobble}

\begin{center}
    \membrete{%
        PONTIFICIA UNIVERSIDAD CATÓLICA DE VALPARAÍSO\\
        FACULTAD DE INGENIERÍA\\
        ESCUELA DE INGENIERÍA INFORMÁTICA\\
    }

    \vspace*{5cm}

    \tituloPortada{%
        Serverless y computación \\
        en el borde
    }

    \vspace*{3.5cm}

    % TODO (equipo): confirmar ortografía/tildes exactas de los 2 nombres
    % marcados con "??" antes de la entrega final.
    \alumnos{%
        Francisca Antonia Abarca Pezo\\
        Vicente André Arratia Caroca\\
        Isidora Valentina Cisternas Cisternas\\
        Rodolfo Antonio Fernández Vera\\
        Valentina Ignacia Guzmán Elgueta\\
        Daniel Ignacio Miranda Vázquez\\
        Matias Ignacio Reyes Pizarro\\
        Fabián Andres Solís Piñones\\
        Claudio Patricio Toledo Mac-lean\\
    }

    \vspace*{1.8cm}

    \datosPortada{%
        \datoPortada{Grupo}{Terabyte}
        \datoPortada{Asignatura}{Taller de Formulación de Proyectos Informáticos}
        \datoPortada{Fecha de entrega}{21 de septiembre de 2026}
    }

    \vspace*{2.0cm}

    \fechaPortada{Septiembre de 2026}
\end{center}

\clearpage


%%%%%%%%%%%%%%%%%%%%%%%%%%%%%%%%%%%%%%%%%%%%%%%%%%%%%%%%%%%%%%%%%%%%%%%%%%%%%%%%
% ÍNDICE GENERAL
%%%%%%%%%%%%%%%%%%%%%%%%%%%%%%%%%%%%%%%%%%%%%%%%%%%%%%%%%%%%%%%%%%%%%%%%%%%%%%%%

\tableofcontents

\clearpage


%%%%%%%%%%%%%%%%%%%%%%%%%%%%%%%%%%%%%%%%%%%%%%%%%%%%%%%%%%%%%%%%%%%%%%%%%%%%%%%%
% LISTA DE FIGURAS Y LISTA DE TABLAS
%%%%%%%%%%%%%%%%%%%%%%%%%%%%%%%%%%%%%%%%%%%%%%%%%%%%%%%%%%%%%%%%%%%%%%%%%%%%%%%%

\phantomsection
\addcontentsline{toc}{section}{Lista de figuras}
\listoffigures

\clearpage
\phantomsection
\addcontentsline{toc}{section}{Lista de tablas}
\listoftables

\clearpage

  %%%%%%%%%%%%%%%%%%%%%%%%%%%%%%%%%%%%%%%%%%%%%%%%%%%%%%%%%%%%%%%%%%%%%%%%%%%%%%%%
% BLOQUE 2 — RESUMEN EJECUTIVO
% Estructura obligatoria del informe TI-05 (punto 2).
% Aquí comienza la numeración arábiga del cuerpo del informe.
%%%%%%%%%%%%%%%%%%%%%%%%%%%%%%%%%%%%%%%%%%%%%%%%%%%%%%%%%%%%%%%%%%%%%%%%%%%%%%%%

\pagenumbering{arabic}
\setcounter{page}{1}

\phantomsection
\addcontentsline{toc}{section}{Resumen ejecutivo}
\section*{Resumen ejecutivo}

% TODO (equipo): redactar el resumen ejecutivo (contexto, alcance,
% principales hallazgos comparativos y recomendación resumida).
% Nota de autoría (punto 6.1 de las indicaciones): esta sección es de
% autoría humana; la IA solo puede usarse para corrección de estilo
% (nivel 1) y debe declararse en el Anexo (bloque 9).

% TODO (opcional): descomentar y completar palabras clave.
% \phantomsection
% \addcontentsline{toc}{subsection}{Palabras clave}
% \palabrasClave{palabra1, palabra2, palabra3}

\clearpage

  %%%%%%%%%%%%%%%%%%%%%%%%%%%%%%%%%%%%%%%%%%%%%%%%%%%%%%%%%%%%%%%%%%%%%%%%%%%%%%%%
% BLOQUE 3 — INTRODUCCIÓN Y CONTEXTO
%%%%%%%%%%%%%%%%%%%%%%%%%%%%%%%%%%%%%%%%%%%%%%%%%%%%%%%%%%%%%%%%%%%%%%%%%%%%%%%%

\phantomsection
\addcontentsline{toc}{section}{Introducción y contexto}
\section*{Introducción y contexto}

% TODO (equipo): presentar el tema TI-05 (serverless y edge computing),
% la empresa asignada (Terabyte), el alcance del trabajo y la
% estructura del informe. Las definiciones conceptuales de base van en
% el Marco teórico (bloque 4), no aquí.
%
% PÁRRAFO OBLIGATORIO DE CIERRE (punto 4 de las indicaciones — NO
% OMITIR): un párrafo breve, al cierre de esta sección, que distinga
% explícitamente qué contenido proviene directamente de la ficha TI-05
% y cuál es aporte propio del grupo (subtemas agregados, alternativas
% adicionales investigadas, mediciones propias, etc.). Este párrafo se
% usa como referencia al evaluar el análisis crítico del grupo.
% Autoría: humana obligatoria (punto 6.1).

\clearpage

  %%%%%%%%%%%%%%%%%%%%%%%%%%%%%%%%%%%%%%%%%%%%%%%%%%%%%%%%%%%%%%%%%%%%%%%%%%%%%%%%
% BLOQUE 4 — MARCO TEÓRICO Y CONCEPTUAL
% Sección agregada tras cruzar la ficha TI-05 con las pautas generales
% del curso, que exigen esta sección entre la Introducción y el
% Desarrollo. Da las bases conceptuales que el Desarrollo (bloque 5)
% profundiza técnicamente.
%%%%%%%%%%%%%%%%%%%%%%%%%%%%%%%%%%%%%%%%%%%%%%%%%%%%%%%%%%%%%%%%%%%%%%%%%%%%%%%%

\section{Marco teórico y conceptual}

% TODO (equipo): definiciones base necesarias para entender el resto
% del informe:
%   - Serverless (arquitectura sin servidor): qué es y qué no es.
%   - FaaS (funciones como servicio).
%   - Contenedores sin servidor.
%   - Edge computing (cómputo en el borde).
%   - El continuo nube-borde (cloud-to-edge continuum): dónde se ubica
%     cada modelo de ejecución entre el centro de datos y el borde de
%     la red.
% Usar fuentes preferentes (documentación oficial, estándares, papers
% revisados por pares) según el punto 5 de las indicaciones del curso.

\clearpage

  %%%%%%%%%%%%%%%%%%%%%%%%%%%%%%%%%%%%%%%%%%%%%%%%%%%%%%%%%%%%%%%%%%%%%%%%%%%%%%%%
% BLOQUE 5 — DESARROLLO Y ANÁLISIS TÉCNICO
% Estructura obligatoria del informe TI-05 (ver ficha, cruzada con las
% pautas generales del curso). Cubre los 6 subtemas obligatorios de la
% ficha TI-05, todos en ESTE ÚNICO ARCHIVO (por preferencia del equipo,
% sin subcarpetas ni \input separados por subtema).
%%%%%%%%%%%%%%%%%%%%%%%%%%%%%%%%%%%%%%%%%%%%%%%%%%%%%%%%%%%%%%%%%%%%%%%%%%%%%%%%

\section{Desarrollo y análisis técnico}

% --- Subtema 1/6 (ficha TI-05) ---
\subsection{Modelos de ejecución y aislamiento}

% TODO (equipo): cubrir funciones como servicio (FaaS), contenedores sin
% servidor y entornos de borde; y los mecanismos de aislamiento V8,
% microVM y contenedor.

\subsubsection{Alternativas adicionales identificadas por el grupo}

% OBLIGATORIO (punto 4 de las indicaciones): la ficha lista, por
% categoría, algunos proveedores/productos y cierra con "y otros que el
% grupo debe identificar". Para cada categoría comparada aquí (FaaS,
% contenedores sin servidor, entornos de borde), agregar al menos DOS
% alternativas adicionales no mencionadas en la ficha, con su
% justificación; o, si tras la búsqueda no se encuentran más,
% declararlo expresamente y documentar la búsqueda (fuentes
% consultadas, criterios de descarte, fecha).

% --- Subtema 2/6 (ficha TI-05) ---
\subsection{Arranque en frío}

% TODO (equipo): causas del arranque en frío, cómo se mide, y
% mitigaciones (concurrencia aprovisionada, instantáneas/snapshots,
% entornos livianos). Ver también el Entregable específico de PoC/
% mediciones de arranque en frío (bloque 6).

% --- Subtema 3/6 (ficha TI-05) ---
\subsection{Facturación por consumo y punto de equilibrio}

% TODO (equipo): unidades de facturación (GB-segundo, vCPU-segundo,
% tiempo de CPU frente a tiempo de reloj, número de invocaciones,
% salida de datos) y el concepto de punto de equilibrio frente a
% instancias reservadas. El cálculo numérico del punto de equilibrio va
% en el Entregable específico "Modelo de costos" (bloque 6); esta
% subsección explica los conceptos y variables que ese modelo usa.

% --- Subtema 4/6 (ficha TI-05) ---
\subsection{Diseño orientado a eventos}

% TODO (equipo): colas, publicación/suscripción, idempotencia, entrega
% al menos una vez ("at-least-once"), y orquestación frente a
% coreografía.

% --- Subtema 5/6 (ficha TI-05) ---
\subsection{WebAssembly y WASI como entorno de borde}

% TODO (equipo): portabilidad, aislamiento y arranque submilisegundo de
% WebAssembly/WASI como entorno de ejecución en el borde.

% --- Subtema 6/6 (ficha TI-05) ---
\subsection{Límites técnicos y dependencia del proveedor (vendor lock-in)}

% TODO (equipo): tiempos máximos de ejecución, tamaño de carga útil,
% manejo de estado y estrategias de portabilidad entre proveedores. Ver
% también el cuadro comparativo de plataformas (bloque 6).

\clearpage

  %%%%%%%%%%%%%%%%%%%%%%%%%%%%%%%%%%%%%%%%%%%%%%%%%%%%%%%%%%%%%%%%%%%%%%%%%%%%%%%%
% BLOQUE 6 — ENTREGABLES ESPECÍFICOS
% Los 3 entregables específicos de la ficha TI-05, integrados en el
% cuerpo del informe (no como anexo). Todos en ESTE ÚNICO ARCHIVO (por
% preferencia del equipo, sin subcarpetas ni \input separados por
% entregable).
%%%%%%%%%%%%%%%%%%%%%%%%%%%%%%%%%%%%%%%%%%%%%%%%%%%%%%%%%%%%%%%%%%%%%%%%%%%%%%%%

\section{Entregables específicos}

% --- Entregable 1/3 (ficha TI-05) ---
\subsection{Modelo de costos: serverless vs. instancias siempre encendidas}

% TODO (equipo): modelo de costos que compare una arquitectura
% serverless con una de contenedores/instancias siempre encendidas,
% identificando el volumen de equilibrio en solicitudes por mes
% (fórmula, tabla o gráfico, según convenga).
%
% Recordatorio (punto 5 de las indicaciones sobre precios y fuentes):
% toda cifra debe indicar producto, moneda, región, fecha de consulta y
% si es precio de lista o cotización; verificar en la calculadora/
% página oficial del proveedor (no reutilizar cifras de terceros sin
% verificar). Si un proveedor no publica precios, declararlo
% expresamente ("solo por cotización").

% --- Entregable 2/3 (ficha TI-05) ---
\subsection{Prueba de concepto y mediciones de arranque en frío / latencia}

% TODO (equipo): elegir UNA de las dos opciones que exige la ficha:
%   (a) Prueba de concepto propia (código + mediciones), documentando
%       el entorno, el método de medición y los resultados; o
%   (b) Análisis de mediciones publicadas de arranque en frío/latencia,
%       citando la fuente y explicando su metodología.
% Si se incluye código propio, usar el entorno `minted` ya cargado por
% la clase (\begin{minted}{lenguaje} ... \end{minted}) y compilar con
% -shell-escape en Overleaf (Menu > Settings, o ya configurado según el
% motor del proyecto).

% --- Entregable 3/3 (ficha TI-05) ---
\subsection{Cuadro comparativo de plataformas}

% TODO (equipo): cuadro comparativo de AL MENOS 6 plataformas con sus
% límites técnicos (tiempo máximo de ejecución, tamaño máximo de
% payload, manejo de estado, portabilidad, etc.). Debe incluir
% plataformas de más de una categoría de la ficha (FaaS, contenedores
% sin servidor, cómputo en el borde, código abierto, datos sin
% servidor).
%
% OBLIGATORIO (punto 4 de las indicaciones): incorporar al menos DOS
% plataformas adicionales no mencionadas en la ficha (de ahí las 8
% filas de la tabla), con su justificación; o declarar expresamente que
% no se encontraron más alternativas relevantes y documentar la
% búsqueda (fuentes, criterios de descarte, fecha).
%
% Filas de la tabla: completar (mínimo 6 + 2 alternativas). Borrar las
% filas sobrantes o agregar más según corresponda.

\tabla{Cuadro comparativo de plataformas serverless y de borde}{cuadro-comparativo-plataformas}{|p{2.6cm}|p{2.3cm}|p{2cm}|p{2cm}|p{2.2cm}|p{3.4cm}|}{%
    \filaTabla{\textbf{Plataforma} & \textbf{Categoría} & \textbf{Tiempo máx. ejecución} & \textbf{Tamaño máx. payload} & \textbf{Manejo de estado} & \textbf{Otros límites / notas}}
    \filaTabla{ & & & & & }
    \filaTabla{ & & & & & }
    \filaTabla{ & & & & & }
    \filaTabla{ & & & & & }
    \filaTabla{ & & & & & }
    \filaTabla{ & & & & & }
    \filaTabla{ & & & & & }
    \filaTabla{ & & & & & }
}

\clearpage

  %%%%%%%%%%%%%%%%%%%%%%%%%%%%%%%%%%%%%%%%%%%%%%%%%%%%%%%%%%%%%%%%%%%%%%%%%%%%%%%%
% BLOQUE 7 — VÍNCULO CON LA PROPUESTA TÉCNICO-ECONÓMICA
% El enunciado deja a criterio del grupo si esto va dentro de
% Desarrollo o de Conclusiones; se dejó como sección propia (entre
% Entregables y Tendencias) para que el orden de la pauta sea
% explícito. Si el equipo prefiere fusionarlo con Desarrollo o
% Conclusiones, basta con mover este \input en main.tex y ajustar el
% \section por \subsection.
%%%%%%%%%%%%%%%%%%%%%%%%%%%%%%%%%%%%%%%%%%%%%%%%%%%%%%%%%%%%%%%%%%%%%%%%%%%%%%%%

\section{Vínculo con la propuesta técnico-económica}

% TODO (equipo): este es "el caso más claro de arquitectura que cambia
% la estructura de costos de CAPEX a OPEX puro" (ficha TI-05). Discutir
% ese cambio y el riesgo de "la factura que crece con el éxito del
% sistema" (a mayor tráfico/adopción, mayor costo variable, a
% diferencia de una inversión fija ya amortizada). Conectar con el
% modelo de costos del bloque 6.

\clearpage

  %%%%%%%%%%%%%%%%%%%%%%%%%%%%%%%%%%%%%%%%%%%%%%%%%%%%%%%%%%%%%%%%%%%%%%%%%%%%%%%%
% BLOQUE 8 — TENDENCIAS Y EVOLUCIÓN FUTURA
% Sección agregada tras cruzar la ficha TI-05 con las pautas generales
% del curso, que la exigen entre el Vínculo técnico-económico y las
% Conclusiones.
%%%%%%%%%%%%%%%%%%%%%%%%%%%%%%%%%%%%%%%%%%%%%%%%%%%%%%%%%%%%%%%%%%%%%%%%%%%%%%%%

\section{Tendencias y evolución futura}

% TODO (equipo): tendencias emergentes y evolución esperada de
% serverless y edge computing (por ejemplo: WebAssembly como runtime de
% borde consolidándose, edge AI/inferencia en el borde, nuevos modelos
% de facturación, convergencia entre FaaS y contenedores, estándares de
% portabilidad, etc.). Si se citan informes de industria o encuestas,
% declarar explícitamente muestra y metodología (punto 5 de las
% indicaciones: fuentes admisibles con declaración expresa).

\clearpage

  %%%%%%%%%%%%%%%%%%%%%%%%%%%%%%%%%%%%%%%%%%%%%%%%%%%%%%%%%%%%%%%%%%%%%%%%%%%%%%%%
% BLOQUE 9 — CONCLUSIONES Y RECOMENDACIONES
%%%%%%%%%%%%%%%%%%%%%%%%%%%%%%%%%%%%%%%%%%%%%%%%%%%%%%%%%%%%%%%%%%%%%%%%%%%%%%%%

\phantomsection
\addcontentsline{toc}{section}{Conclusiones y recomendaciones}
\section*{Conclusiones y recomendaciones}

% TODO (equipo): conclusiones y una recomendación concreta con la que
% el grupo se comprometa (punto 4 de las indicaciones).
% Autoría: 100% humana, obligatoria (punto 6.1); no se admite uso de IA
% más allá de corrección de estilo (nivel 1), y debe declararse en el
% Anexo (bloque 11).

\clearpage

  %%%%%%%%%%%%%%%%%%%%%%%%%%%%%%%%%%%%%%%%%%%%%%%%%%%%%%%%%%%%%%%%%%%%%%%%%%%%%%%%
% BLOQUE 10 — BIBLIOGRAFÍA / REFERENCIAS
% Las entradas se agregan en referencias.bib (raíz del proyecto) y se
% citan con \parencite{clave} donde corresponda en el cuerpo del
% informe. Estilo: APA (biblatex, backend biber), ya configurado en
% inf-pucv.cls.
%
% Cada entrada debe incluir: autor, año, y según el tipo, volumen,
% número, páginas o DOI. Si la entrada es una cifra de precio, además
% debe indicar fecha de consulta y región (punto 5 de las indicaciones).
% Fuentes preferentes / admisibles / no admisibles: ver punto 5 de las
% indicaciones del curso.
%%%%%%%%%%%%%%%%%%%%%%%%%%%%%%%%%%%%%%%%%%%%%%%%%%%%%%%%%%%%%%%%%%%%%%%%%%%%%%%%

\phantomsection
\addcontentsline{toc}{section}{Bibliografía}

\printbibliography[title={Bibliografía}]

\clearpage

  %%%%%%%%%%%%%%%%%%%%%%%%%%%%%%%%%%%%%%%%%%%%%%%%%%%%%%%%%%%%%%%%%%%%%%%%%%%%%%%%
% BLOQUE 11 — ANEXO OBLIGATORIO: DECLARACIÓN DE USO DE INTELIGENCIA
% ARTIFICIAL
%
% OBLIGATORIO, no opcional: la ficha TI-05 2026 es explícita — "Todo
% informe incluye un anexo de declaración, aunque el nivel declarado
% sea 0 en todas las secciones" y "un informe sin anexo de declaración
% no se recibe conforme" (punto 6.3 de las indicaciones). Esta ficha
% específica prevalece sobre una diapositiva más genérica de pautas
% generales que lo listaba como opcional.
%
% Formulario base en el punto 6.3 y en el Anexo de las indicaciones del
% curso, adaptado aquí con el nivel declarado por cada una de las 10
% secciones REALES de este informe (bloques 1 a 10; ya no las 4
% genéricas del formulario base) y con una tabla de firmas de 9 filas
% (una por integrante).
%
% Reinicia la numeración de página con prefijo "A" (A1, A2, ...), igual
% que en la plantilla original.
%
% Cada integrante declara su propio uso; nadie declara por otro.
%%%%%%%%%%%%%%%%%%%%%%%%%%%%%%%%%%%%%%%%%%%%%%%%%%%%%%%%%%%%%%%%%%%%%%%%%%%%%%%%

\pagenumbering{arabic}
\setcounter{page}{1}
\renewcommand{\thepage}{\hspace*{1em}A\arabic{page}}

\phantomsection
\addcontentsline{toc}{section}{Anexo: Declaración de uso de inteligencia artificial}
\section*{Anexo: Declaración de uso de inteligencia artificial}

\textoPortadaFont{\textbf{Empresa:} Terabyte\\}
\textoPortadaFont{\textbf{Tema:} TI-05 — Serverless y computación en el borde (edge computing)\\}
% TODO (equipo): completar fecha de la declaración.
\textoPortadaFont{\textbf{Fecha:} \\}

\bigskip

\subsection*{A. Declaración por sección}

% Nivel 0: sin uso de IA. Nivel 1: asistencia superficial (estilo/
% corrección). Nivel 2: colaboración e ideación. Nivel 3: generación
% sustancial. Ver punto 6.2 de las indicaciones para las definiciones
% completas. Los niveles 2 y 3 requieren evidencia trazable en la
% columna D / sección C de este anexo.
%
% TODO (equipo): esta tabla lista cada una de las 10 secciones reales
% de este informe (bloques 1 a 10). Completar nivel, herramienta/
% versión y evidencia para cada fila.

{\scriptsize
\begin{longtable}{|p{5.6cm}|p{1.6cm}|p{3.4cm}|p{4.3cm}|}
    \hline
    \textbf{Sección del informe} & \textbf{Nivel (0-3)} & \textbf{Herramienta y versión} & \textbf{Evidencia adjunta (niveles 2 y 3)} \\
    \hline
    \endhead
    1. Portada y formalidades & & & \\ \hline
    2. Resumen ejecutivo & & & \\ \hline
    3. Introducción y contexto & & & \\ \hline
    4. Marco teórico y conceptual & & & \\ \hline
    5. Desarrollo y análisis técnico & & & \\ \hline
    6. Entregables específicos & & & \\ \hline
    7. Vínculo con la propuesta técnico-económica & & & \\ \hline
    8. Tendencias y evolución futura & & & \\ \hline
    9. Conclusiones y recomendaciones & & & \\ \hline
    10. Bibliografía & & & \\ \hline
\end{longtable}
}

\subsection*{B. Prompts utilizados (niveles 2 y 3)}

% TODO (equipo): transcribir aquí, o en hoja/enlace adjunto
% identificado, los prompts efectivamente empleados, indicando a qué
% sección corresponde cada uno.

\subsection*{C. Evidencia trazable adjunta (niveles 2 y 3)}

% TODO (equipo): indicar el medio por el cual se entrega la evidencia
% (enlace a la conversación, archivo exportado, historial de versiones
% del documento u Overleaf, o repositorio con sus commits) y verificar
% que el profesor pueda acceder a ella.

\subsection*{D. Firma de los integrantes}

% "Declaro que la información consignada en este formulario es
% completa y veraz, que las secciones señaladas en el punto 6.1 de las
% indicaciones son de autoría humana y que puedo explicar oralmente
% cualquier parte del trabajo entregado."
%
% TODO (equipo): confirmar ortografía/tildes de los nombres marcados
% con "??" (ver también portada, bloque 1) antes de imprimir y firmar
% esta página. Cada integrante completa su propia fila (nivel máximo
% declarado y firma); nadie firma por otro.

{\scriptsize
\begin{longtable}{|p{0.8cm}|p{6.2cm}|p{3.8cm}|p{4.1cm}|}
    \hline
    \textbf{N.\textordmasculine} & \textbf{Nombre completo} & \textbf{Nivel máximo declarado} & \textbf{Firma} \\
    \hline
    \endhead
    1 & Francisca Antonia Abarca Pezo & & \\ \hline
    2 & Vicente André Arratia Caroca & & \\ \hline
    3 & Isidora Valentina Cisternas Cisternas & & \\ \hline
    4 & Rodolfo Antonio Fernández Vera & & \\ \hline
    5 & Valentina Ignacia Guzmán Elgueta & & \\ \hline
    6 & Daniel Ignacio Miranda Vázquez & & \\ \hline
    7 & Matias Ignacio Reyes Pizarro & & \\ \hline
    8 & Fabián Andres Solís Piñones & & \\ \hline
    9 & Claudio Patricio Toledo Mac-lean & & \\ \hline
\end{longtable}
}

\clearpage

  %%%%%%%%%%%%%%%%%%%%%%%%%%%%%%%%%%%%%%%%%%%%%%%%%%%%%%%%%%%%%%%%%%%%%%%%%%%%%%%%
% BLOQUE 12 — DOCUMENTO ADJUNTO: CUESTIONARIO DE 30 PREGUNTAS
% Requisitos en el punto 3 de las indicaciones del curso.
%
% Reinicia la numeración de página con prefijo "B" (B1, B2, ...) para
% distinguirla del Anexo de declaración de IA (prefijo "A", bloque 11).
%%%%%%%%%%%%%%%%%%%%%%%%%%%%%%%%%%%%%%%%%%%%%%%%%%%%%%%%%%%%%%%%%%%%%%%%%%%%%%%%

\pagenumbering{arabic}
\setcounter{page}{1}
\renewcommand{\thepage}{\hspace*{1em}B\arabic{page}}

\phantomsection
\addcontentsline{toc}{section}{Documento adjunto: Cuestionario de evaluación de conocimientos}
\section*{Documento adjunto: Cuestionario de evaluación de conocimientos}

% TODO (equipo) — requisitos obligatorios del cuestionario (punto 3 de
% las indicaciones):
%   - 30 preguntas en total sobre aspectos clave de las tecnologías
%     investigadas.
%   - Combinación de formatos: selección múltiple, verdadero/falso,
%     completar y respuesta corta.
%   - Cada pregunta incluye su respuesta correcta y una breve
%     justificación.
%   - Cubrir tanto teoría como aplicaciones prácticas.
%   - Dificultad graduada: 40% básicas, 40% intermedias, 20% avanzadas.
%   - Incluir preguntas que evalúen comprensión, aplicación y análisis
%     crítico (no solo memorización).
%   - Autoría: la redacción de las preguntas y justificaciones es
%     obligatoriamente humana (punto 6.1); no se admite generación por
%     IA en esta sección.
%
% Sugerencia de organización: agrupar las preguntas por sección de
% origen (ver índice temático abajo) y marcar el nivel de dificultad de
% cada una junto a su enunciado.

\subsection*{Índice temático}

% TODO (equipo): completar cuántas preguntas corresponden a cada una de
% las 10 secciones del informe y su rango (ej. "1-3"), de modo que
% sumen 30.

{\scriptsize
\begin{longtable}{|p{7.5cm}|p{2.5cm}|p{2.5cm}|}
    \hline
    \textbf{Sección del informe} & \textbf{N.\textordmasculine\ de preguntas} & \textbf{Rango} \\
    \hline
    \endhead
    1. Portada y formalidades & & \\ \hline
    2. Resumen ejecutivo & & \\ \hline
    3. Introducción y contexto & & \\ \hline
    4. Marco teórico y conceptual & & \\ \hline
    5. Desarrollo y análisis técnico & & \\ \hline
    6. Entregables específicos & & \\ \hline
    7. Vínculo con la propuesta técnico-económica & & \\ \hline
    8. Tendencias y evolución futura & & \\ \hline
    9. Conclusiones y recomendaciones & & \\ \hline
    10. Bibliografía & & \\ \hline
    \textbf{Total} & & 30 \\ \hline
\end{longtable}
}

\subsection*{Preguntas}

% TODO (equipo): redactar aquí las 30 preguntas con su respuesta y
% justificación. No hay una macro predefinida para esto en la clase;
% se sugiere un formato simple, por ejemplo:
%
%   \textbf{P1 (básica — Introducción y contexto).} Enunciado de la
%   pregunta. \textit{Respuesta:} X. \textit{Justificación:} breve
%   explicación.

\clearpage


\end{document}
